\chapter{Reinforcement Learning with Verifiable Rewards: GRPO}\label{ch:grpo}

Chapters~\ref{ch:scaling} and~\ref{ch:refine} improved reasoning by spending more
compute \emph{per query}: sample many answers and vote, or let the model critique and
revise itself. Every one of those tricks leaves the weights untouched, so the bill is paid
again on every request. This chapter takes the opposite trade. We spend compute
\emph{once}, during training, so that the model itself becomes a better reasoner and every
later query is cheaper. The method is reinforcement learning (RL) with \emph{verifiable
rewards} (RLVR), and the optimiser is Group Relative Policy Optimization (GRPO), the
algorithm behind DeepSeek-R1 and many of the open reasoning models that followed it.

For a systems person, RL post-training is the most interesting workload in the LLM
lifecycle because it fuses the two workloads we have studied separately. Each step is an
\emph{inference} job (generate thousands of tokens per prompt, autoregressively,
memory-bound) followed by a \emph{training} job (a differentiable forward and backward
pass over the same tokens, compute-bound). Get the balance between the two wrong and
either your GPUs idle while the sampler crawls, or your sampler idles while the optimiser
waits. In this chapter we build GRPO from first principles, keeping the code small enough
to read in one sitting, and we mark where each design decision has a systems cost.
\Cref{ch:grpo2} then makes the algorithm robust, and \cref{ch:rlsys} scales it into a
distributed system.

\begin{figure}[t]
\centering
\resizebox{\linewidth}{!}{%
\begin{tikzpicture}[
  st/.style={draw=sysblue,rounded corners=3pt,fill=sysblue!8,minimum height=11mm,
    text width=24mm,align=center,font=\small},
  im/.style={st,draw=labpurple,fill=labpurple!8},
  num/.style={circle,fill=#1,text=white,font=\scriptsize\bfseries,inner sep=1.5pt},
  >={Stealth[length=2.2mm]}]
\node[st] (a) at (0,0) {RL for LLMs\\(\S\ref{sec:grpo:rl})};
\node[st] (b) at (3,0) {RLHF\\recap};
\node[st] (c) at (6,0) {Verifiable\\rewards};
\node[st] (d) at (9,0) {GRPO at a\\glance};
\node[im] (e) at (9,-2.1) {Model \&\\MATH data};
\node[im] (f) at (6,-2.1) {Rollouts $\to$ rewards\\$\to$ advantages};
\node[im] (g) at (3,-2.1) {Log-probs $\to$\\loss $\to$ loop};
\node[im] (h) at (0,-2.1) {Checkpoints \&\\evaluation};
\foreach \n/\c/\i in {a/sysblue/1,b/sysblue/2,c/sysblue/3,d/sysblue/4,e/labpurple/5,f/labpurple/6,g/labpurple/7,h/labpurple/8}
  \node[num=\c] at (\n.north west) {\i};
\draw[->,thick,sysblue!70] (a)--(b); \draw[->,thick,sysblue!70] (b)--(c);
\draw[->,thick,sysblue!70] (c)--(d); \draw[->,thick,labpurple!70] (d)--(e);
\draw[->,thick,labpurple!70] (e)--(f); \draw[->,thick,labpurple!70] (f)--(g);
\draw[->,thick,labpurple!70] (g)--(h);
\node[font=\footnotesize\itshape,sysblue,anchor=south] at (4.5,0.75) {concepts};
\node[font=\footnotesize\itshape,labpurple,anchor=north] at (4.5,-2.85) {implementation (Lab 27)};
\end{tikzpicture}}
\caption{Roadmap of this chapter. The top row builds the conceptual picture (why RL, how
RLHF works, what changes when rewards are verifiable, what GRPO computes). The bottom
row implements GRPO stage by stage and trains a model with it.}
\label{fig:grpo:roadmap}
\end{figure}

% ======================================================================
\section{Reinforcement learning for LLMs}\label{sec:grpo:rl}

There are two independent dials for buying reasoning accuracy with compute
(\cref{fig:grpo:dials}). \emph{Inference-time scaling} spends more FLOPs on each answer:
longer chains of thought, more samples, more refinement rounds. \emph{Training-time
scaling} spends FLOPs up front to change the weights. The two compose: a model improved by
RL can still be wrapped in self-consistency voting at inference time, and in practice the
strongest systems do both.

\begin{figure}[ht]
\centering
\begin{tikzpicture}[>={Stealth[length=2mm]},font=\small]
\begin{scope}
\draw[->] (0,0)--(5.2,0) node[below left,font=\footnotesize]{queries served};
\draw[->] (0,0)--(0,3.2) node[above,font=\footnotesize,align=center]{cumulative\\compute};
\draw[thick,warnred] (0,0.15)--(5,3.0) node[right,font=\footnotesize,align=left]{inference-time\\scaling};
\draw[thick,keygreen] (0,1.3)--(5,2.0) node[right,font=\footnotesize,align=left]{training-time\\scaling};
\draw[thick,sysblue,dashed] (0,0.15)--(5,1.15) node[right,font=\footnotesize]{baseline};
\draw[decorate,decoration={brace,amplitude=3pt},keygreen] (-0.1,0.15)--(-0.1,1.3)
  node[midway,left,font=\scriptsize,align=right]{one-off\\RL cost};
\end{scope}
\end{tikzpicture}
\caption{Two ways to buy accuracy. Inference-time scaling raises the marginal cost of
every query; training-time scaling pays a fixed cost once and keeps the per-query cost
close to the baseline. Which is cheaper depends on how many queries the model will serve.}
\label{fig:grpo:dials}
\end{figure}

RL in the LLM world is almost always a \emph{post-training} stage. A pretrained model has
absorbed knowledge from next-token prediction; RL changes \emph{how} it uses that
knowledge, by optimising a property of the whole output (is the final answer right? do
humans prefer it?) rather than the likelihood of each individual token. Two RL stages are
common (\cref{fig:grpo:pipeline}): \emph{reasoning training} with verifiable rewards,
which teaches the model to reach correct answers on math and code, and \emph{preference
tuning}, which nudges style and helpfulness towards what people like. The order varies
between labs and the two are sometimes interleaved. DeepSeek-R1-Zero famously skipped
both instruction tuning and preference tuning and ran reasoning RL directly on the base
model. The result is less polished than a full pipeline, but it shows clean reasoning
behaviour and, for us, has a big methodological advantage: any improvement can be
attributed to RL alone.

\begin{figure}[ht]
\centering
\resizebox{\linewidth}{!}{%
\begin{tikzpicture}[box/.style={draw,rounded corners=2pt,minimum height=9mm,text width=23mm,
   align=center,font=\small},>={Stealth[length=2mm]}]
\node[box,fill=fdata!15] (pt) {Pretraining\\{\scriptsize next-token, $\sim10^{13}$ tok}};
\node[box,fill=fdata!15,right=7mm of pt] (sft) {Instruction\\fine-tuning (SFT)};
\node[box,fill=fcompute!15,right=7mm of sft] (rlvr) {Reasoning RL\\{\scriptsize RLVR: correctness}};
\node[box,fill=fcomm!15,right=7mm of rlvr] (rlhf) {Preference RL\\{\scriptsize RLHF: human taste}};
\draw[->,thick] (pt)--(sft); \draw[->,thick] (sft)--(rlvr); \draw[->,thick] (rlvr)--(rlhf);
\draw[->,thick,dashed,warnred] (pt.south) to[bend right=18]
  node[below,font=\scriptsize,warnred]{R1-Zero shortcut: RLVR directly on the base model} (rlvr.south);
\node[font=\scriptsize,align=center,above=1mm of sft]{Ch.~\ref{ch:instruct}};
\node[font=\scriptsize,align=center,above=1mm of rlvr]{this chapter};
\node[font=\scriptsize,align=center,above=1mm of pt]{Ch.~\ref{ch:pretrain}};
\end{tikzpicture}}
\caption{Where RL sits in an LLM training pipeline. The order of the two RL stages is not
fixed; some pipelines interleave them. The dashed path applies reasoning RL directly to a
pretrained base model, which is the setting of this chapter.}
\label{fig:grpo:pipeline}
\end{figure}

\paragraph{The RL vocabulary, translated.} Classic RL talks about agents, environments,
states and actions. For a language model the mapping is straightforward
(\cref{tab:grpo:vocab}). The unusual features of the LLM setting are that the action space
is huge ($|\mathcal{V}|\approx\num{150000}$ tokens per step), the state is the entire prefix,
transitions are deterministic (appending a token), and the reward is \emph{sparse}: a
single number at the end of a trajectory that may be thousands of tokens long.

\begin{table}[ht]
\centering\small
\begin{tabularx}{\linewidth}{@{}lX@{}}
\toprule
RL term & Meaning for an LLM\\
\midrule
policy $\pitheta$ & the language model with weights $\theta$\\
state $s_t$ & prompt $x$ plus the tokens generated so far, $y_{<t}$\\
action $a_t$ & the next token $y_t\in\mathcal{V}$\\
trajectory / rollout / completion & one full sampled response $y=(y_1,\dots,y_{|y|})$\\
reward $R(x,y)$ & a scalar for the whole response (here: 1 if the final answer is right)\\
episode & one prompt $\to$ one completion; ends at EOS or the token budget\\
on-policy & the data were sampled from the current weights\\
\bottomrule
\end{tabularx}
\caption{RL terminology mapped onto autoregressive generation. A ``rollout'' is the
generation process and a ``completion'' is its output; the words are used interchangeably.}
\label{tab:grpo:vocab}
\end{table}

\paragraph{The policy-gradient theorem in one line.} (\Cref{ch:pg} develops the foundations in full.) We want to maximise the expected
reward $J(\theta)=\E_{x\sim\mathcal{D}}\,\E_{y\sim\pitheta(\cdot\mid x)}[R(x,y)]$. The
reward is a non-differentiable function of discrete samples, so we cannot backpropagate
through it. The \emph{log-derivative trick}
$\nabla_\theta \pitheta(y\mid x)=\pitheta(y\mid x)\nabla_\theta\log\pitheta(y\mid x)$ moves
the gradient onto the log-probability, which \emph{is} differentiable:
\begin{equation}
\nabla_\theta J(\theta)
 = \E_{x,\,y\sim\pitheta}\!\big[\,R(x,y)\,\nabla_\theta\log\pitheta(y\mid x)\big]
 = \E\Big[R(x,y)\sum_{t=1}^{|y|}\nabla_\theta\log\pitheta(y_t\mid x,y_{<t})\Big].
\label{eq:grpo:pg}
\end{equation}
The second equality is the autoregressive factorisation
$\log\pitheta(y\mid x)=\sum_t\log\pitheta(y_t\mid x,y_{<t})$. A Monte Carlo estimate of
\cref{eq:grpo:pg} needs only samples and their log-probabilities: sample, score, and push
up the log-probability of each sample in proportion to its reward. This is the REINFORCE
estimator (Williams, 1992). It is unbiased but very noisy. The standard fix is to subtract a
\emph{baseline} $b(x)$ that does not depend on $y$:
\begin{equation}
\E_{y\sim\pitheta}\big[b(x)\nabla_\theta\log\pitheta(y\mid x)\big]
 = b(x)\sum_y\nabla_\theta\pitheta(y\mid x)
 = b(x)\nabla_\theta\underbrace{\textstyle\sum_y\pitheta(y\mid x)}_{=1}=0,
\end{equation}
so replacing $R$ by the \emph{advantage} $A=R-b(x)$ leaves the gradient unbiased and can
reduce its variance dramatically. Every algorithm in this chapter and the next is a choice
of (i) how to estimate the baseline, (ii) how to weight and clip each token's contribution,
and (iii) how far the policy may move per update.

\subsection{The original pipeline: RL from human feedback}\label{sec:grpo:rlhf}

RL entered mainstream LLM training through InstructGPT (Ouyang et al., 2022),
the recipe that turned GPT-3 into the first ChatGPT. Its RL from human feedback (RLHF)
pipeline has two phases (\cref{fig:grpo:rlhf}).

\emph{Phase 1: learn a reward model.} For many prompts, sample several responses (with the
temperature and nucleus sampling of \cref{ch:scaling}) and ask annotators to rank them.
Take a question like ``Which laptop specs matter if I mostly program?'' One answer lists
fifteen components; another says ``prioritise RAM and an SSD, and a GPU only if you train
models locally''. Annotators tend to prefer the specific answer. A \emph{reward model}
$r_\phi(x,y)$, usually an LLM with a scalar head, is fitted to these comparisons with the
Bradley--Terry likelihood (\cref{ch:eval}):
\begin{equation}
\mathcal{L}_{\text{RM}}(\phi)=-\E_{(x,y^+,y^-)}\log\sigma\big(r_\phi(x,y^+)-r_\phi(x,y^-)\big).
\end{equation}

\emph{Phase 2: optimise the policy against the reward model.} Generate responses, score
them with $r_\phi$, and update the policy with Proximal Policy Optimization (PPO)
(Schulman et al., 2017). PPO maximises a clipped surrogate (\cref{ch:grpo2})
using per-token advantages from a learned \emph{value model} $V_\psi(s_t)$ (the critic)
via generalised advantage estimation, and it subtracts a per-token KL penalty
$\beta\log\frac{\pitheta}{\piref}$ from the reward so the policy does not drift into text
that fools the reward model.

\begin{figure}[ht]
\centering
\resizebox{\linewidth}{!}{%
\begin{tikzpicture}[b/.style={draw,rounded corners=2pt,align=center,font=\small,minimum height=8mm},
  >={Stealth[length=2mm]}]
\node[font=\small\bfseries,anchor=west] at (-0.3,1.1) {Phase 1: fit a reward model};
\node[b,fill=fdata!12] (p1) at (0.6,0) {prompt};
\node[b,fill=fcompute!12] (l1) at (2.6,0) {policy\\LLM};
\node[b,fill=white,text width=24mm] (r1) at (5.4,0) {responses\\A, B, C, D};
\node[b,fill=fcomm!12] (h1) at (8.4,0) {humans rank\\{\scriptsize C $\succ$ A $\succ$ D $\succ$ B}};
\node[b,fill=napkin!15] (rm) at (12.5,0) {reward model\\$r_\phi(x,y)\in\R$};
\draw[->] (p1)--(l1); \draw[->] (l1)--(r1); \draw[->] (r1)--(h1);
\draw[->] (h1)--node[above,font=\scriptsize]{BT loss}(rm);
\node[font=\small\bfseries,anchor=west] at (-0.3,-1.3) {Phase 2: optimise the policy};
\node[b,fill=fdata!12] (p2) at (0.6,-2.4) {prompt};
\node[b,fill=fcompute!12] (l2) at (2.6,-2.4) {policy\\$\pitheta$};
\node[b,fill=white] (r2) at (5.4,-2.4) {response $y$};
\node[b,fill=napkin!15] (rm2) at (8.4,-2.4) {$r_\phi(x,y)$};
\node[b,fill=fcompute!25,text width=27mm] (ppo) at (12.5,-2.4) {PPO update\\{\scriptsize + value model, KL to $\piref$}};
\draw[->] (p2)--(l2); \draw[->] (l2)--(r2); \draw[->] (r2)--(rm2); \draw[->] (rm2)--(ppo);
\draw[->,thick,fcompute] (ppo.south) -- ++(0,-0.5) -| node[pos=0.25,below,font=\scriptsize]{gradient step on $\theta$} (l2.south);
\draw[->,dashed,napkin] (rm.south)--(rm2.north);
\end{tikzpicture}}
\caption{RLHF. Phase 1 turns human rankings into a learned scalar reward. Phase 2 samples
from the policy, scores each sample with the frozen reward model and updates the policy
with PPO. The same skeleton (generate, score, update) survives in RLVR; only the scorer
changes.}
\label{fig:grpo:rlhf}
\end{figure}

RLHF matters to us as the ancestor of RLVR. The control flow is the one we will implement:
generate, score, update. What makes RLHF heavy is the cast of models. A PPO-based RLHF step
touches four networks of comparable size: the policy (trained), the value model (trained),
the reward model (frozen) and the reference policy (frozen).

\begin{napkinmath}[the four-model tax of PPO-RLHF]
Take an 8B policy with mixed-precision AdamW at 16 bytes per trained parameter
(\cref{ch:pretrain}).
\begin{itemize}
\item Policy (trained): $8\times10^9\times16\,\text{B}=128$\,GB.
\item Value model (trained, same size): another 128\,GB.
\item Reward model and reference model (frozen, bf16): $2\times 8\times10^9\times2\,\text{B}=32$\,GB.
\end{itemize}
That is $\approx288$\,GB of state before a single activation or KV-cache byte, so the job
spans at least four 80\,GB GPUs just to hold it. GRPO drops the value model (no critic),
and RLVR drops the reward model (a program replaces it). What remains is the policy
plus, optionally, the reference: $\approx128+16=144$\,GB, roughly half.
\end{napkinmath}

\subsection{From human feedback to verifiable rewards}\label{sec:grpo:rlvr}

Training a reward model is expensive, and a learned scorer is also \emph{exploitable}:
push hard enough and the policy finds outputs the reward model likes but humans do not
(reward hacking). For tasks where correctness can be \emph{checked}, both problems
disappear. RL with verifiable rewards (RLVR) replaces $r_\phi$ by a deterministic program
(\cref{fig:grpo:rlvr}). For math, that program is the answer verifier of \cref{ch:eval}:
extract the final \verb|\boxed{}| answer, normalise it, test equivalence with the reference,
and return 1 or 0. For code it is a compiler plus unit tests run in a sandbox.

\begin{figure}[ht]
\centering
\begin{tikzpicture}[b/.style={draw,rounded corners=2pt,align=center,font=\small,minimum height=8mm},
  >={Stealth[length=2mm]}]
\node[b,fill=fdata!12,text width=27mm] (q) {``Compute $7\cdot(2+4)-9$. Put the answer in \texttt{\textbackslash boxed\{\}}.''};
\node[b,fill=fcompute!12,right=6mm of q] (pi) {policy $\pitheta$};
\node[b,text width=27mm,right=6mm of pi] (y) {``$2+4=6$, $7\cdot 6=42$, $42-9=33$.
 \texttt{\textbackslash boxed\{33\}}''};
\node[b,fill=keygreen!15,right=6mm of y] (v) {verifier\\{\scriptsize vs.\ reference ``33''}};
\node[b,fill=keygreen!35,below=7mm of v] (r) {$R=1$};
\draw[->] (q)--(pi); \draw[->] (pi)--(y); \draw[->] (y)--(v); \draw[->] (v)--(r);
\draw[->,thick,fcompute] (r.west) -| node[pos=0.3,below,font=\scriptsize]{policy-gradient update} (pi.south);
\end{tikzpicture}
\caption{RLVR collapses RLHF's two phases into one loop. A program (here, an answer
checker) produces the reward, so there is no reward model to train, store or serve.}
\label{fig:grpo:rlvr}
\end{figure}

RLVR became the default recipe for reasoning after DeepSeek-R1 (DeepSeek-AI, 2025)
, which showed that correctness checks on math plus compile-and-test checks on code,
with no human preference data and no learned reward model, are enough to grow long,
self-correcting chains of thought. Its advantages:
\begin{itemize}
\item \textbf{No reward model.} One fewer network to train, store and run. The verifier
is CPU code that costs microseconds to milliseconds per sample.
\item \textbf{Determinism.} The same output always gets the same reward. There is no
annotator disagreement and no reward-model noise.
\item \textbf{Unlimited labels.} Given reference answers, any number of fresh samples can
be scored without more human effort.
\end{itemize}
The price is \emph{scope}: RLVR only applies where a checker exists (math, code, formal
proofs, structured extraction, games). It is also only as good as the checker. A verifier
that accepts \verb|\boxed{33.0000001}| or that reads the \emph{first} box when the model
writes ten will be gamed. Code rewards add an operational burden: model-generated programs
must run in an isolated sandbox with CPU, memory and time limits and no network access, a
small distributed system of its own (\cref{ch:rlsys}). The RL algorithm is identical either
way, so we stay with math.

% ======================================================================
\section{RLVR with GRPO}\label{sec:grpo:grpo}

RLVR fixes \emph{where the learning signal comes from}. We still need a policy optimiser
that turns rewards into weight updates. PPO is the classic choice. GRPO, introduced in
DeepSeekMath (Shao et al., 2024) and used for DeepSeek-R1, is a simplification
that removes PPO's value model. Instead of learning $V_\psi(s)$ to predict how good a state
is, GRPO samples a \emph{group} of $G$ completions for the same prompt and uses the group's
own reward statistics as the baseline. That is the ``group relative'' in the name.

\begin{figure}[ht]
\centering
\resizebox{0.95\linewidth}{!}{%
\begin{tikzpicture}[m/.style={draw,rounded corners=2pt,minimum width=20mm,minimum height=8mm,
  align=center,font=\small},>={Stealth[length=2mm]}]
\node[font=\bfseries] at (2.4,1.4) {PPO};
\node[m,fill=fcompute!20] (pp) at (0,0) {policy\\{\scriptsize trained}};
\node[m,fill=fcompute!20] (pv) at (2.4,0) {value $V_\psi$\\{\scriptsize trained}};
\node[m,fill=gray!15] (pr) at (4.8,0) {reference\\{\scriptsize frozen}};
\node[m,fill=gray!15] (prm) at (2.4,-1.3) {reward model\\{\scriptsize frozen}};
\node[font=\scriptsize,align=center] at (2.4,-2.35) {advantage $A_t$ from GAE over $V_\psi$\\one sample per prompt is enough};
\begin{scope}[xshift=8cm]
\node[font=\bfseries] at (2.4,1.4) {GRPO + RLVR};
\node[m,fill=fcompute!20] (gp) at (0.8,0) {policy\\{\scriptsize trained}};
\node[m,fill=gray!15] (gr) at (4.0,0) {reference\\{\scriptsize frozen, optional}};
\node[m,fill=keygreen!15] (gv) at (2.4,-1.3) {verifier\\{\scriptsize a program}};
\node[font=\scriptsize,align=center] at (2.4,-2.35) {advantage $A_i$ from the group of $G$\\samples of the \emph{same} prompt};
\end{scope}
\draw[dashed,gray] (6.9,1.6)--(6.9,-2.7);
\end{tikzpicture}}
\caption{The cast of models. PPO-based RLHF trains a critic and runs a reward model. GRPO
with verifiable rewards replaces the critic by group statistics and the reward model by a
program. The only extra network left is an optional frozen reference for a KL penalty.}
\label{fig:grpo:cast}
\end{figure}

GRPO is not free, though: it moves cost from \emph{memory} (no critic) to \emph{sampling}
(G completions per prompt instead of one). Since sampling is the expensive phase of RL
(\cref{sec:grpo:loop}), that trade is worth understanding precisely, and much of
\cref{ch:rlsys} is about making it cheaper.

\subsection{An intuition: practising a climbing route}\label{sec:grpo:analogy}

GRPO has many moving parts and a lot of RL jargon, so before the equations here is a
picture to hang them on. Imagine a climber training on an indoor bouldering wall.
\begin{enumerate}
\item The coach points at a \textbf{route} (the \emph{prompt}).
\item The climber makes several \textbf{attempts} at it, each a little different because
their movement is not perfectly repeatable (the \emph{rollouts}, sampled at temperature
$>0$).
\item The only feedback is binary and comes at the very end: did they \textbf{top out}?
Nobody grades individual moves (a \emph{sparse, outcome reward}).
\item After the session the climber compares attempts \textbf{with each other}: an
attempt that topped out on a route they usually fail is very informative, while an
attempt that topped out on a route they always send is not (the \emph{group-relative
advantage}).
\item For every attempt they recall how \textbf{habitual} its moves were, that is, how
likely they would be to repeat that exact sequence (the \emph{log-probability}).
\item They adjust their habits: make the move sequences of better-than-average attempts
more habitual and the others less so (the \emph{policy-gradient update}).
\item Optionally they check they are not drifting away from the fundamentals their
first coach drilled into them (the \emph{KL penalty to a reference policy}) and avoid
overhauling their style in a single session (\emph{clipping}, \cref{ch:grpo2}).
\end{enumerate}
Notice what is \emph{absent}: no one trained a separate judge to predict how good an
attempt will be halfway up the wall. That judge would be PPO's value model, and GRPO does
without it by comparing whole attempts on the same route.

\subsection{The GRPO procedure}\label{sec:grpo:procedure}

Formally, for a prompt $x$ GRPO samples a group $\{y_1,\dots,y_G\}\sim\piold(\cdot\mid x)$,
scores each with the verifier to get $r_1,\dots,r_G$, and forms the advantages
\begin{equation}
A_i=\frac{r_i-\operatorname{mean}(r_{1:G})}{\operatorname{std}(r_{1:G})+\epsilon}.
\label{eq:grpo:adv}
\end{equation}
Every token of completion $i$ shares the same advantage $A_i$. The full objective of
DeepSeekMath, which we build up to over two chapters, is
\begin{equation}
\mathcal{J}(\theta)=\E\Bigg[\frac1G\sum_{i=1}^{G}\frac1{|y_i|}\sum_{t=1}^{|y_i|}
\Big(\min\big(\rho_{i,t}A_i,\ \clip(\rho_{i,t},1-\varepsilon,1+\varepsilon)A_i\big)
-\beta\,\mathbb{D}_{i,t}\Big)\Bigg],
\label{eq:grpo:full}
\end{equation}
with the per-token importance ratio
$\rho_{i,t}=\pitheta(y_{i,t}\mid x,y_{i,<t})/\piold(y_{i,t}\mid x,y_{i,<t})$, and where $\mathbb{D}_{i,t}$ is a per-token estimate of $\KL(\pitheta\,\|\,\piref)$. That is a
lot at once. This chapter implements the core, the advantage-weighted log-likelihood
\begin{equation}
\mathcal{L}_{\text{PG}}(\theta)=-\frac1G\sum_{i=1}^G A_i\,\log\pitheta(y_i\mid x)
 =-\frac1G\sum_{i=1}^G A_i\sum_{t}\log\pitheta(y_{i,t}\mid x,y_{i,<t}),
\label{eq:grpo:pg_loss}
\end{equation}
which is REINFORCE with a group-mean baseline and group-std scaling. \Cref{ch:grpo2} adds
the ratio and clipping, the KL term and format rewards, and examines the many variants
(token- vs.\ sequence-level averaging, no std scaling, \dots) that recent work proposes.
\Cref{fig:grpo:step} shows the six stages for one prompt with numbers we compute later in
this chapter. Keep it at hand; every section below fills in one box.

% (figure fig:grpo:step is defined in the worked-example section below)

% ======================================================================
\section{Loading a pretrained model}\label{sec:grpo:load}

GRPO needs a policy that \emph{sometimes} succeeds. If every rollout gets reward 0, every
advantage in \cref{eq:grpo:adv} is 0 and nothing is learned (\cref{sec:grpo:adv}). That is
why reasoning RL starts from a capable pretrained model rather than from random weights.
At full scale we start from the Qwen3-0.6B \emph{base} checkpoint of \cref{ch:modern},
loaded with the same helper used throughout Part~VI:

\begin{lstlisting}[caption={Loading the base policy and its tokenizer (real-hardware path of \texttt{run.py}).},label={lst:grpo:load}]
from llmsys.qwen3 import load_qwen3_from_dir, Qwen3Tokenizer
from llmsys.utils import get_device

device = get_device()                       # cuda > mps > cpu
model = load_qwen3_from_dir("Qwen3-0.6B-Base", "0.6B",
                            dtype=torch.bfloat16, device=device)
tok = Qwen3Tokenizer("Qwen3-0.6B-Base/tokenizer.json")
tok.eos_token_id = tok.pad_token_id         # base model ends text with <|endoftext|>
\end{lstlisting}

Before training, sample one answer with the temperature and top-$p$ sampler of
\cref{ch:scaling} to make sure everything is wired correctly. A base model asked
``Half of $3x-9$ equals $x+37$; find $x$'' will often produce a wrong or unboxed answer,
which is fine: we only need the plumbing to work.

\paragraph{A CPU-sized stand-in.} A single GRPO step on Qwen3-0.6B generates tens of
thousands of tokens, and the experiment is only practical on a GPU. To make every result in
this chapter reproducible on a laptop we also use a \emph{toy verifiable task}
(\code{llmsys.rl.toy}): single-digit addition. The prompt is \verb|Q:3+5=\nA:| and a
fully rewarded completion is \verb|\boxed{8}| followed by \verb|<|\verb+|eos|>+. A 2-layer,
64-wide Qwen3-shaped model (about \num{76000} parameters, with the same RMSNorm, RoPE, GQA
and SwiGLU code as the real one) is first \emph{SFT-warm-started} on deliberately
imperfect demonstrations: 40\% of answers are off by one, only 30\% are boxed, and 20\%
open with a \verb|<think>...</think>| scratchpad. This mimics the situation of a real base
model: it ``knows'' the arithmetic most of the time but rarely presents the answer in the
format the verifier insists on. Every curve in \cref{fig:grpo:curves} and every number in
\cref{fig:grpo:step} comes from this toy, measured on a CPU in this book's build
environment.

\begin{pitfall}[\code{inference\_mode} versus \code{no\_grad}]
The generation helpers of earlier chapters were wrapped in
\code{@torch.inference\_mode()}. Tensors created inside inference mode are flagged and can
\emph{never} participate in autograd later, so reusing the sampled ids or cached
activations in the differentiable log-probability pass raises an error. RL code that
samples and then trains must sample under \code{@torch.no\_grad()}, which disables
gradient tracking without tagging the outputs.
\end{pitfall}

% ======================================================================
\section{Loading a MATH training subset}\label{sec:grpo:data}

The MATH dataset (Hendrycks et al., 2021) has about \num{12500} competition problems.
Five hundred of them form MATH-500, the evaluation set of \cref{ch:eval}. For training we
use the remaining \num{12000}, a split built to exclude every MATH-500 problem
(\cref{fig:grpo:split}), so that evaluation numbers measure generalisation rather than
memorisation. Contamination is a systems problem as much as a statistical one: once
training data is assembled from many sources by many scripts, someone must
\emph{prove} that the test set is not in it (\cref{ch:eval} discusses $n$-gram overlap
checks).

\begin{figure}[ht]
\centering
\begin{tikzpicture}[font=\small]
\fill[fdata!35] (0,0) rectangle (9.6,0.7);
\fill[warnred!45] (9.6,0) rectangle (10,0.7);
\draw (0,0) rectangle (10,0.7);
\node at (4.8,0.35) {training split: \num{12000} problems (this chapter, Ch.~\ref{ch:distill})};
\node[anchor=west,font=\scriptsize] at (10.1,0.35) {MATH-500};
\draw[decorate,decoration={brace,amplitude=4pt,mirror}] (0,-0.1)--(10,-0.1)
  node[midway,below=4pt,font=\scriptsize]{MATH: $\approx$\num{12500} problems, 7 subjects, difficulty levels 1--5};
\end{tikzpicture}
\caption{Disjoint training and evaluation splits of MATH. RL samples only from the left
part; all reported accuracies use the right part.}
\label{fig:grpo:split}
\end{figure}

A record has the fields \code{problem}, \code{answer}, \code{solution}, \code{type}
(subject), \code{level} and an id:

\begin{lstlisting}[style=plain,caption={Shape of one training record (contents abbreviated).},label={lst:grpo:record}]
{"problem":  "A rental costs $12 per day plus $0.15 per mile. ... How many miles?",
 "answer":   "180",
 "solution": "Let m be the number of miles ... so m = \\boxed{180}.",
 "type": "Algebra", "level": "Level 2", "unique_id": 917}
\end{lstlisting}

RLVR uses only \code{problem} (the prompt) and \code{answer} (what the verifier compares
against). The worked \code{solution} is deliberately \emph{not} used. Rewarding similarity
to a reference derivation would push the model towards one particular style of proof,
whereas an outcome reward lets it explore any path that ends at the right answer. This is
the essential difference from supervised fine-tuning on solutions (\cref{ch:distill}).
For the toy task, \code{make\_examples} enumerates the 100 facts $a+b$ with
$a,b\in\{0,\dots,9\}$. Because the toy's ``knowledge'' comes from SFT on all 100 facts, we
train and evaluate on the same prompts and read evaluation as \emph{in-distribution}
accuracy; the code also supports a held-out split for the exercises.

% ======================================================================
\section{Sampling rollouts}\label{sec:grpo:rollouts}

Stage 1 of GRPO samples $G$ completions per prompt from the current policy.
\Cref{lst:grpo:sample} is the whole sampler: prefill the prompt into a KV cache
(\cref{ch:inference}), then decode one token at a time with temperature and top-$p$
(\cref{ch:scaling}) until EOS or the token budget.

\begin{lstlisting}[caption={Unbatched rollout sampler (\texttt{labs/lab17\_grpo/solution.py}).},label={lst:grpo:sample}]
@torch.no_grad()
def sample_response(model, tokenizer, prompt, max_new_tokens=32,
                    temperature=1.0, top_p=1.0, generator=None):
    device = next(model.parameters()).device
    prompt_ids = torch.tensor(tokenizer.encode(prompt), device=device)
    cache = KVCache(model.n_layers)
    logits = model(prompt_ids[None], cache=cache)[:, -1]   # prefill
    generated = []
    for _ in range(max_new_tokens):                        # decode
        nxt = sample_next_token(logits, temperature, None, top_p,
                                generator)
        generated.append(int(nxt))
        if int(nxt) == tokenizer.eos_token_id:
            break
        logits = model(nxt, cache=cache)[:, -1]
    full = torch.cat([prompt_ids,
                      torch.tensor(generated, device=device)])
    return full, prompt_ids.numel(), tokenizer.decode(generated)
\end{lstlisting}

Three details matter. First, the function returns the \emph{full} id sequence and the
prompt length, which is exactly what the log-probability stage needs, so nothing has to be
re-tokenised (re-tokenising decoded text is a classic RL bug because
$\text{encode}(\text{decode}(y))\neq y$ for many BPE sequences). Second, the EOS token is
\emph{kept} in the completion. It is a real action the policy took, and if it were never
trained on, a long RL run could slowly forget how to stop. Third, it uses an explicit
\code{torch.Generator} so a rollout can be replayed bit for bit in tests.

Calling \code{sample\_response} $G$ times in a loop is the clearest form, and the version the
chapter's walkthrough uses. It is also the slowest. The library sampler
\code{llmsys.rl.rollout.sample\_rollouts} stacks all $P\times G$ sequences of a step into one
left-padded batch (\code{generate\_batch} from \cref{ch:serving}). Since decode is
memory-bound, a batch of 64 costs barely more per step than a batch of 1.

\begin{systemslens}[a GRPO step is mostly an inference job]
Consider $P\times G=64$ rollouts of up to 1024 new tokens after a 256-token prompt.
Using the roofline helper \code{llmsys.flops.decode\_step\_time\_lower\_bound} on an H100:
\begin{itemize}
\item \textbf{Qwen3-0.6B.} Every decode step must stream 1.19\,GB of weights plus a KV
cache that reaches 9.4\,GB at 1280 tokens (64 sequences). Summed over 1024 steps, the
bandwidth bound alone is $\approx$\,2.1\,s. The training side (old-logprob forward plus
forward/backward on $64\times1280$ tokens, $\approx6N\!\cdot\!\text{tokens}=2.9\times10^{14}$
FLOPs) takes $\approx$\,0.3\,s at peak, or 0.74\,s at a realistic 40\% MFU.
\item \textbf{Qwen3-8B.} Rollouts: $\approx$\,7.2\,s lower bound. Training:
$4.0\times10^{15}$ FLOPs, $\approx$\,10\,s at 40\% MFU.
\end{itemize}
So ``rollouts dominate'' is not a law of nature. It depends on model size, response length
and how close the sampler gets to the bandwidth roofline. Real samplers rarely exceed 50\%
of peak bandwidth, reasoning responses are 4--32$\times$ longer than 1024 tokens, and the
\emph{slowest} rollout in the batch sets the step time (the long-tail problem of
\cref{ch:rlsys}). In practice, rollouts take most of the wall-clock time in reasoning RL,
which is why production systems run them on a dedicated inference engine (vLLM, SGLang)
rather than in the trainer.
\end{systemslens}

% ======================================================================
\section{Calculating rewards}\label{sec:grpo:rewards}

Stage 2 scores each rollout. Our reward is the verifier of \cref{ch:eval} in its strictest
mode: the answer must appear inside \verb|\boxed{}|, and the boxed content must match the
reference after normalisation.

\begin{lstlisting}[caption={The RLVR reward: correct \emph{and} boxed.},label={lst:grpo:reward}]
def reward_rlvr(text, truth):
    pred = extract_boxed(text)          # last \boxed{...}, nested braces ok
    if pred is None:
        return 0.0                      # implicit format requirement
    return float(normalize_answer(pred) == normalize_answer(truth))
\end{lstlisting}

The \code{None} branch is doing quiet but important work. A response that says ``the answer
is 8'' gets nothing, even though it is right. The reward therefore teaches two things at
once, correctness and the output contract that the evaluation harness relies on. For the
four toy completions of the prompt \verb|Q:4+3=| shown in \cref{fig:grpo:step} the rewards
are $(1,1,0,0)$: the third is correct but unboxed, the fourth is boxed but wrong.

\begin{keyidea}[outcome rewards beat process rewards (so far)]
It is tempting to reward intermediate steps with a \emph{process reward model} (PRM).
The DeepSeek-R1 report describes trying this and abandoning it: per-step rewards were hard
to define, expensive to compute and easy to hack. Rewarding only the final answer and
letting the model discover its own intermediate steps worked better. Process supervision
remains useful at \emph{inference} time for reranking (\cref{ch:refine}).
\end{keyidea}

\begin{pitfall}[verifier loopholes become policy behaviour]
Whatever the verifier tolerates, RL will find. A checker that reads the \emph{first}
\verb|\boxed{}| teaches the model to write many boxes; one that does substring matching
teaches it to box long lists of numbers. Extract the \emph{last} box, require it to be
unique if you can afford to, cap the length of the boxed content, and unit-test the
verifier against adversarial strings before any long run.
\end{pitfall}

% ======================================================================
\section{From rewards to advantages}\label{sec:grpo:adv}

A reward says how good one rollout is in absolute terms. An \emph{advantage} says how much
better or worse it is than its siblings for the same prompt (\cref{eq:grpo:adv}):

\begin{lstlisting}[caption={Group-relative advantages.},label={lst:grpo:adv}]
def compute_advantages(rewards, eps=1e-4):
    rewards = rewards.float()
    return (rewards - rewards.mean()) / (rewards.std() + eps)   # unbiased std
\end{lstlisting}

For rewards $(1,1,0,0)$ the mean is $0.5$, the unbiased standard deviation is
$\sqrt{1/3}=0.577$, and the advantages are $(+0.866,+0.866,-0.866,-0.866)$. The advantage
multiplies the gradient of each rollout's log-likelihood: positive values make that
rollout more likely, negative values less likely, and values near zero contribute almost
nothing.

\paragraph{Binary rewards in closed form.} With $r_i\in\{0,1\}$ and a fraction $p$ of the
group correct, $\operatorname{mean}=p$ and the unbiased std is
$s=\sqrt{\tfrac{G}{G-1}p(1-p)}$. A correct rollout therefore gets
$A^+=(1-p)/s$ and an incorrect one $A^-=-p/s$. \Cref{fig:grpo:advplot} plots both for
$G=8$. When a prompt is hard (small $p$), the rare success gets a large positive push and
each failure a small negative one; when it is easy, the reverse. Because the advantages sum
to zero, the total ``probability mass'' moved is balanced within every group.

\begin{figure}[ht]
\centering
\begin{tikzpicture}
\begin{axis}[width=0.72\linewidth,height=5.2cm,xlabel={fraction of the group that is correct, $p$},
  ylabel={advantage},xmin=0,xmax=1,ymin=-3,ymax=3,grid=major,grid style={gray!20},
  legend style={font=\scriptsize,at={(0.98,0.97)},anchor=north east},
  xtick={0,0.125,0.25,0.5,0.75,0.875,1},xticklabels={0,$\frac18$,$\frac14$,$\frac12$,$\frac34$,$\frac78$,1},
  tick label style={font=\scriptsize}]
\addplot[very thick,keygreen,domain=0.125:0.875,samples=60] {(1-x)/sqrt(8/7*x*(1-x))};
\addlegendentry{$A^+$ (correct rollout)}
\addplot[very thick,warnred,domain=0.125:0.875,samples=60] {-x/sqrt(8/7*x*(1-x))};
\addlegendentry{$A^-$ (wrong rollout)}
\addplot[only marks,mark=*,keygreen,mark size=1.6pt] coordinates {(0.125,2.475)(0.25,1.620)(0.375,1.208)(0.5,0.935)(0.625,0.725)(0.75,0.540)(0.875,0.354)};
\addplot[only marks,mark=*,warnred,mark size=1.6pt] coordinates {(0.125,-0.354)(0.25,-0.540)(0.375,-0.725)(0.5,-0.935)(0.625,-1.208)(0.75,-1.620)(0.875,-2.475)};
\addplot[only marks,mark=square*,gray,mark size=2pt] coordinates {(0,0)(1,0)};
\node[font=\scriptsize,gray,anchor=south west] at (axis cs:0.01,0.15) {all wrong: $A\equiv0$};
\node[font=\scriptsize,gray,anchor=north east] at (axis cs:0.99,-0.2) {all right: $A\equiv0$};
\end{axis}
\end{tikzpicture}
\caption{Advantages for binary rewards with $G=8$ (dots are the attainable values
$p=k/8$). Rare successes on hard prompts receive the largest positive weight. At $p=0$ and
$p=1$ every advantage is zero and the prompt contributes no gradient.}
\label{fig:grpo:advplot}
\end{figure}

\paragraph{Zero-variance groups.} If all $G$ rollouts get the same reward, every advantage
is zero (the $\epsilon$ in the denominator prevents a division by zero) and the policy
gradient is exactly zero. This is the right behaviour: a group that is uniformly right or
uniformly wrong says nothing about \emph{which} behaviour to reinforce. Its cost is real,
though. We paid for $G$ full rollouts and got no gradient. With per-prompt success rate $p$
the probability of a zero-variance group is $p^G+(1-p)^G$, which for $G=8$ is 27\% at
$p=0.15$ and 43\% at $p=0.9$. \Cref{ch:grpo2} tracks this fraction and discusses DAPO's
dynamic sampling, which spends extra rollouts to replace such groups.

\begin{keyidea}[GRPO's baseline is a Monte Carlo value estimate]
PPO learns $V_\psi(x)$, the expected reward from prompt $x$. GRPO estimates the same
quantity by averaging $G$ samples, $\hat V(x)=\frac1G\sum_j r_j$. It is unbiased, needs no
training, and its variance shrinks as $1/G$. It is \emph{not} independent of $y_i$, since
$r_i$ is part of its own baseline, which introduces a small $O(1/G)$ bias; leave-one-out
baselines (RLOO) remove it. Dividing by the group std is a separate heuristic that
\cref{ch:grpo2} revisits (Dr.~GRPO drops it).
\end{keyidea}

% ======================================================================
\section{Scoring rollouts with sequence log-probabilities}\label{sec:grpo:logprobs}

Stage 4 asks how likely the current policy was to produce each rollout. In
\cref{ch:refine} we scored answers by their \emph{average} per-token log-probability,
\begin{equation}
\bar\ell(y\mid x)=\frac1{|y|}\sum_{t=1}^{|y|}\log\pitheta(y_t\mid x,y_{<t}),
\end{equation}
because length normalisation makes answers of different lengths comparable. GRPO, as we
implement it in this chapter, uses the \emph{sum}, $\log\pitheta(y\mid x)=|y|\cdot\bar\ell$.
The reward and the advantage belong to the whole sequence, so the quantity they multiply in
\cref{eq:grpo:pg_loss} should be the log-likelihood of the whole sequence. Averaging would
silently reweight the update: a 2000-token rollout and a 20-token rollout with the same
advantage would contribute equal total gradient, so each token of the long one would be
pushed 100 times less. Whether that is desirable is a real design question.
\Cref{ch:grpo2} shows that the original GRPO paper actually \emph{does} average per sequence,
and that DAPO and Dr.~GRPO each criticise a different consequence of that choice. For now,
the summed version is the clean REINFORCE estimator.

A side effect of summing is that sequence log-probabilities scale with length: each extra
token adds a negative term. For two equally rewarded answers, the shorter one has the
larger log-likelihood, and gradient ascent on $A\cdot\log\pi$ is ``cheaper'' along it. In the
worked example (\cref{fig:grpo:step}) the 11-token answer with a \verb|<think>| scratchpad has
log-probability $-3.06$, versus $-2.27$ for the bare 4-token boxed answer.

\begin{lstlisting}[caption={Differentiable sequence log-probability. \texttt{gather} picks, for every position, the log-prob of the token that actually came next.},label={lst:grpo:seqlp}]
def sequence_logprob(model, token_ids, prompt_len):
    logits = model(token_ids[None])[0].float()   # (T, V), upcast to fp32
    logprobs = torch.log_softmax(logits, dim=-1)
    # row t predicts token t+1: pick the log-prob of the realised next token
    picked = logprobs[:-1].gather(1, token_ids[1:, None]).squeeze(-1)
    return picked[prompt_len - 1:].sum()         # completion tokens only
\end{lstlisting}

The indexing deserves a careful look (\cref{fig:grpo:shift}). Row $t$ of the logits
predicts token $t+1$, so the first completion token (index \code{prompt\_len}) is scored by
row \code{prompt\_len-1}. Slicing the \code{(T-1,)} vector from \code{prompt\_len-1} therefore
keeps exactly the $|y|$ completion terms. A first draft might build an explicit
\code{torch.arange} of positions and index with it. \code{gather} computes the same thing in
one fused indexing kernel and avoids materialising index tensors.

\begin{figure}[ht]
\centering
\resizebox{0.95\linewidth}{!}{%
\begin{tikzpicture}[c/.style={draw,minimum width=9mm,minimum height=6mm,font=\scriptsize\ttfamily},
  >={Stealth[length=1.6mm]}]
\foreach \t/\s/\f in {0/Q/fdata,1/:/fdata,2/4/fdata,3/+/fdata,4/3/fdata,5/=/fdata,6/\textbackslash n/fdata,7/A/fdata,8/:/fdata,9/\textbackslash box/fcompute,10/7/fcompute,11/\}/fcompute,12/eos/fcompute}
  \node[c,fill=\f!15] (i\t) at (\t*0.95,0) {\s};
\node[anchor=east,font=\scriptsize] at (-0.6,0) {ids};
\foreach \t in {0,...,11} \node[c,fill=gray!8] (l\t) at (\t*0.95,-1.1) {$\ell_{\t}$};
\node[anchor=east,font=\scriptsize] at (-0.6,-1.1) {logits row};
\foreach \t in {8,...,11} {\pgfmathtruncatemacro{\n}{\t+1}\draw[->,fcompute,thick] (l\t.north) -- (i\n.south);}
\draw[decorate,decoration={brace,amplitude=3pt,mirror},fcompute] (8*0.95-0.45,-1.5)--(11*0.95+0.45,-1.5)
 node[midway,below=3pt,font=\scriptsize,fcompute]{rows \texttt{prompt\_len-1 .. T-2} score the completion};
\draw[decorate,decoration={brace,amplitude=3pt},fdata] (-0.45,0.4)--(8*0.95+0.45,0.4)
 node[midway,above=3pt,font=\scriptsize,fdata!80!black]{prompt, \texttt{prompt\_len}=9};
\end{tikzpicture}}
\caption{The off-by-one in log-probability extraction. Logit row $t$ predicts token $t+1$,
so the completion (blue) is scored by rows $8,\dots,11$ for a 9-token prompt. The
prompt's own likelihood is ignored, since the policy did not choose those tokens.}
\label{fig:grpo:shift}
\end{figure}

The returned scalar carries a \code{grad\_fn}. PyTorch has recorded the graph from the
weights through the embedding, every attention and MLP block, the LM head and the
\code{log\_softmax} to this sum, so a later \code{backward()} will deliver
$\nabla_\theta\log\pitheta(y\mid x)$, which is exactly the term the policy-gradient theorem asks for.

\begin{systemslens}[the hidden $(T\times V)$ tensor]
\code{log\_softmax(logits.float())} materialises a float32 tensor of shape
$(N, T, |\mathcal V|)$ for $N$ rollouts. For Qwen3 ($|\mathcal V|=\num{151936}$), 8 rollouts
of 2048 tokens need $8\times2048\times\num{151936}\times4\,\text{B}\approx10$\,GB, plus the
same again for its gradient in the backward pass. That is often \emph{larger than the
model}. Production trainers therefore (i) compute log-probs in row chunks, (ii) use fused
``linear + cross-entropy'' kernels that never materialise the full logits (for example
Liger-Kernel's fused linear cross-entropy), and (iii) recompute rather than store. Note
also that we \emph{recompute} log-probs with a fresh forward pass instead of keeping the
logits from sampling. Keeping them would pin $T$ autograd graphs from the decode loop,
and in a real system the sampler is a different engine whose numerics differ from the
trainer's anyway (\cref{ch:rlsys}).
\end{systemslens}

% ======================================================================
\section{From advantages to a policy update: the GRPO loss}\label{sec:grpo:loss}

Stage 5 combines the two branches. Multiply each rollout's sequence log-probability by its
advantage, average over the group, and negate (\cref{eq:grpo:pg_loss}):

\begin{lstlisting}[caption={The policy-gradient loss.},label={lst:grpo:pgloss}]
def policy_gradient_loss(advantages, logps):
    return -(advantages.detach() * logps).mean()
\end{lstlisting}

The negation is only a convention. We want to \emph{maximise}
$\sum_iA_i\log\pitheta(y_i\mid x)$, and PyTorch optimisers \emph{minimise}, so we minimise the
negative. Lowering $-f$ from $-2$ to $-3$ is the same as raising $f$ from 2 to 3.
\code{.detach()} marks the advantages as constants. They depend on rewards, which are not
differentiable, and in a more elaborate setup (a learned reward or a baseline network) we
would not want gradients leaking into them.

Differentiating makes the mechanics explicit:
\begin{equation}
\nabla_\theta\mathcal{L}_{\text{PG}}=-\frac1G\sum_{i=1}^G A_i\sum_t\nabla_\theta\log\pitheta(y_{i,t}\mid x,y_{i,<t}).
\end{equation}
A gradient-descent step moves $\theta$ along $+A_i\nabla\log\pi(y_i)$. Every token of a
positive-advantage rollout becomes more probable, and every token of a negative one less
probable, \emph{including tokens the two rollouts share}. When a correct and an incorrect
rollout both begin with \verb|\boxed{|, the shared prefix receives opposite pushes that
partly cancel. What survives is the difference: the \verb|7| versus the \verb|8|.
Credit assignment in GRPO emerges statistically from such cancellations across many groups,
not from any per-token judgement.

\begin{figure}[t]
\centering
\resizebox{\linewidth}{!}{%
\begin{tikzpicture}[font=\small,
  cell/.style={draw=gray!60,minimum height=7mm,anchor=center},
  hdr/.style={font=\scriptsize\bfseries,align=center},
  stage/.style={circle,fill=#1,text=white,font=\scriptsize\bfseries,inner sep=1.4pt},
  >={Stealth[length=2mm]}]
\node[draw=fdata,fill=fdata!10,rounded corners=2pt,align=center,font=\small] (pr) at (6.2,2.0)
  {prompt $x$: \texttt{Q:4+3=}\textbackslash\texttt{n A:} \quad reference answer: 7};
\node[stage=fdata,anchor=east] at (pr.west) {1};
\node[draw=fcompute,fill=fcompute!12,rounded corners=2pt] (pol) at (6.2,1.05) {policy $\pitheta$ samples $G=4$ rollouts};
\draw[->] (pr)--(pol);
% header
\node[hdr] at (0.6,0.2) {advantage\\$A_i$};
\node[hdr] at (2.2,0.2) {reward\\$r_i$};
\node[hdr] at (6.2,0.2) {rollout $y_i$};
\node[hdr] at (10.0,0.2) {$\log\pitheta(y_i\mid x)$};
\node[hdr] at (11.9,0.2) {$|y_i|$};
\foreach \y/\a/\r/\txt/\lp/\L/\col in {
  -0.6/+0.866/1/{\textbackslash boxed\{7\}}/-2.269/4/keygreen,
  -1.35/+0.866/1/{<think>4+3=7</think>\textbackslash boxed\{7\}}/-3.059/11/keygreen,
  -2.1/-0.866/0/{7 \ (correct, not boxed)}/-1.299/2/warnred,
  -2.85/-0.866/0/{\textbackslash boxed\{8\}}/-3.400/4/warnred} {
  \node[cell,fill=\col!15,minimum width=12mm] at (0.6,\y) {$\a$};
  \node[cell,fill=\col!15,minimum width=10mm] at (2.2,\y) {$\r$};
  \node[cell,fill=white,minimum width=58mm,font=\scriptsize\ttfamily] at (6.2,\y) {\txt};
  \node[cell,fill=fcompute!8,minimum width=18mm] at (10.0,\y) {$\lp$};
  \node[cell,fill=white,minimum width=9mm] at (11.9,\y) {$\L$};
}
\node[stage=keygreen] at (2.2,0.75) {2};
\node[stage=keygreen] at (0.6,0.75) {3};
\node[stage=fcompute] at (10.0,0.75) {4};
\node[draw=sysblue,fill=sysblue!10,rounded corners=2pt,align=center] (loss) at (6.2,-4.0)
  {$\mathcal{L}_{\text{PG}}=-\frac14\sum_iA_i\log\pitheta(y_i\mid x)=+0.136$};
\node[stage=sysblue,anchor=east] at (loss.west) {5};
\draw[->] (0.6,-3.2) |- (loss.west);
\draw[->] (10.0,-3.2) |- (loss.east);
\node[draw=sysblue,fill=sysblue!25,rounded corners=2pt] (upd) at (6.2,-5.0) {backward $+$ optimiser step on $\theta$};
\node[stage=sysblue,anchor=east] at (upd.west) {6};
\draw[->] (loss.south)--(upd.north);
\end{tikzpicture}}
\caption{One GRPO step on the toy task, with numbers computed from the SFT-warm-started
model. (1) One prompt; (2) verifier rewards; (3) group-relative advantages; (4) summed
log-probabilities from a differentiable forward pass; (5) the loss; (6) the update. The
most likely rollout under the current policy, the unboxed ``7'', is penalised, because the
reward requires the box.}
\label{fig:grpo:step}
\end{figure}

\Cref{fig:grpo:step} traces the stages on our toy prompt. The loss is
$-\tfrac14[0.866(-2.269)+0.866(-3.059)-0.866(-1.299)-0.866(-3.400)]=+0.136$. Its value is
\emph{not} a measure of quality. Unlike a supervised loss it can be positive or negative,
it does not trend towards zero, and two runs with identical learning can show very
different loss curves. Only its gradient matters. Note what the update will do here: the
unboxed ``7'' is currently the policy's favourite completion ($-1.30$ is the highest
log-probability in the group), and it will become less likely, while both boxed correct
answers become more likely. That is how RLVR teaches the output contract.

\begin{deepdive}[why the summed loss equals ``ratio $\times$ advantage'' at ratio 1]
\Cref{ch:grpo2} writes the objective with the importance ratio
$\rho_i=\pitheta(y_i)/\piold(y_i)$ instead of $\log\pitheta(y_i)$. Since
$\nabla_\theta\rho_i=\rho_i\nabla_\theta\log\pitheta(y_i)$ and $\rho_i=1$ when
$\theta=\theta_{\text{old}}$, the two objectives have identical gradients on the first
update after sampling. They differ only when the same rollouts are reused for several
updates, which is exactly when clipping becomes necessary.
\end{deepdive}

% ======================================================================
\section{Putting the stages into one function}\label{sec:grpo:onefn}

\Cref{lst:grpo:compute} packages stages 1--5 for one prompt. It interleaves stages 2 and 4
inside the sampling loop (score each rollout and compute its log-probability as soon as it
exists) to avoid a second loop. The logical dataflow is unchanged.

\begin{lstlisting}[caption={\texttt{compute\_grpo\_loss}: stages 1--5 for one training example.},label={lst:grpo:compute}]
def compute_grpo_loss(model, tokenizer, example, render_prompt,
                      num_rollouts=4, max_new_tokens=16, temperature=1.0,
                      top_p=1.0, generator=None):
    assert num_rollouts >= 2, "GRPO needs a group to compare against"
    prompt = render_prompt(example)
    was_training = model.training
    model.eval()                      # no dropout while sampling/scoring
    logps, rewards, samples = [], [], []
    for _ in range(num_rollouts):
        # stage 1: sample a rollout
        ids, plen, text = sample_response(model, tokenizer, prompt,
                                          max_new_tokens, temperature,
                                          top_p, generator)
        r = reward_rlvr(text, example["answer"])             # stage 2
        logps.append(sequence_logprob(model, ids, plen))     # stage 4
        rewards.append(r)
        samples.append({"text": text, "reward": r,
                        "gen_len": ids.numel() - plen})
    if was_training:
        model.train()
    adv = compute_advantages(torch.tensor(rewards))          # stage 3
    loss = policy_gradient_loss(adv, torch.stack(logps))     # stage 5
    return {"loss": loss.item(), "loss_tensor": loss, "rewards": rewards,
            "advantages": adv.tolist(), "samples": samples}
\end{lstlisting}

The function returns both the Python float (for logging) and the tensor (for
\code{backward()}), plus per-sample diagnostics. Run it on a prompt the policy cannot solve
and you will see rewards $[0,0,\dots]$, advantages $[0,0,\dots]$ and a loss of exactly zero:
a wasted step, which is the zero-variance case of \cref{sec:grpo:adv}.

\begin{systemslens}[memory: why the loop above does not scale]
Each \code{sequence\_logprob} call keeps its autograd graph alive until \code{backward()}.
With $G$ rollouts that means $G$ sets of saved activations at once, $G$ times the memory
of one training sequence. At $G=8$ and 2048 tokens on Qwen3-0.6B this is several GB of
activations before counting the $(T\times V)$ logits. The batched trainer
(\code{llmsys.rl.trainer}) fixes this the standard way. It computes the ``old'' log-probs of all
rollouts under \code{no\_grad}, then runs the differentiable pass in \emph{micro-batches}
(\code{minibatches} in \code{GRPOConfig}), calling \code{backward()} after each, so peak
activation memory is that of one micro-batch. Gradient accumulation over micro-batches gives
the same update as one big batch for the summed loss (and the same up to the averaging
constant for the other aggregations of \cref{ch:grpo2}).
\end{systemslens}

% ======================================================================
\section{The GRPO training loop}\label{sec:grpo:loop}

With the loss in hand, the outer loop is an ordinary PyTorch training loop
(\cref{fig:grpo:loop}). Only one box, the loss, is specific to RL.

\begin{figure}[ht]
\centering
\resizebox{\linewidth}{!}{%
\begin{tikzpicture}[b/.style={draw,rounded corners=2pt,minimum height=9mm,text width=22mm,align=center,font=\small},
  n/.style={circle,fill=#1,text=white,font=\scriptsize\bfseries,inner sep=1.3pt},>={Stealth[length=2mm]}]
\node[b,fill=gray!10] (s1) at (0,0) {init policy\\+ AdamW};
\node[b,fill=gray!10] (s2) at (3,0) {next example\\(step $k$)};
\node[b,fill=gray!10] (s3) at (6,0) {\code{zero\_grad()}};
\node[b,fill=fcompute!25,text width=26mm] (s4) at (9.4,0) {\textbf{GRPO loss}\\stages 1--5};
\node[b,fill=gray!10] (s5) at (9.4,-1.9) {\code{backward()}\\+ clip $\|g\|\le1$};
\node[b,fill=gray!10,text width=26mm] (s6) at (6,-1.9) {\code{optimizer.step()}};
\node[b,fill=fdata!15] (s7) at (3,-1.9) {log reward,\\length, loss};
\node[b,fill=fdata!15] (s8) at (0,-1.9) {checkpoint\\every $c$ steps};
\foreach \x/\c/\i in {s1/gray/1,s2/gray/2,s3/gray/3,s4/fcompute/4,s5/gray/5,s6/gray/6,s7/keygreen/7,s8/keygreen/8}
 \node[n=\c] at (\x.north west) {\i};
\draw[->] (s1)--(s2); \draw[->] (s2)--(s3); \draw[->] (s3)--(s4); \draw[->] (s4)--(s5);
\draw[->] (s5)--(s6); \draw[->] (s6)--(s7); \draw[->] (s7)--(s8);
\draw[->,dashed] (s8.north) to[bend left=20] node[left,font=\scriptsize]{repeat} (s2.south west);
\end{tikzpicture}}
\caption{The GRPO training loop. Stages 1--3 and 5--6 are boilerplate; 7--8 are
optional bookkeeping (metrics and checkpoints). Only stage 4, the loss, is RL-specific.}
\label{fig:grpo:loop}
\end{figure}

\begin{lstlisting}[caption={The training loop (abridged from \texttt{train\_rlvr\_grpo}; CSV logging and checkpointing elided).},label={lst:grpo:loop}]
def train_rlvr_grpo(model, tokenizer, data, render_prompt, steps=50,
                    num_rollouts=4, max_new_tokens=16, temperature=1.0,
                    top_p=1.0, lr=3e-4, grad_clip=1.0, seed=0):
    gen = torch.Generator().manual_seed(seed)
    opt = torch.optim.AdamW(model.parameters(), lr=lr, weight_decay=0.0)
    for step in range(steps):
        opt.zero_grad(set_to_none=True)
        ex = data[step % len(data)]
        stats = compute_grpo_loss(model, tokenizer, ex, render_prompt,
                                  num_rollouts, max_new_tokens,
                                  temperature, top_p, gen)
        stats["loss_tensor"].backward()
        torch.nn.utils.clip_grad_norm_(model.parameters(), grad_clip)
        opt.step()
        n = len(stats["samples"])
        reward_avg = sum(stats["rewards"]) / n
        resp_len = sum(x["gen_len"] for x in stats["samples"]) / n
        print(f"[{step+1}/{steps}] loss={stats['loss']:.4f} "
              f"reward_avg={reward_avg:.3f} len={resp_len:.1f}")
\end{lstlisting}

\paragraph{Hyperparameters.} On real models, a temperature of 0.7--1.0 with top-$p$
around 0.9 gives enough diversity within a group for the advantages to be informative,
while keeping completions coherent. The learning rate is small, about $10^{-6}$ to
$10^{-5}$ for a 0.6B model with AdamW, because every step changes a policy that is
simultaneously generating its own training data, and a large step can trigger a feedback
loop of degenerate samples. The group size $G$ (4--16) and the response budget (512 to
tens of thousands of tokens) are mostly a compute decision. Gradient clipping at norm 1 is
cheap insurance against the occasional huge-advantage group.

\paragraph{Reading the output.} Expect noisy logs. With $G=4$ and binary rewards the
per-step reward can only take the values $0,\frac14,\frac12,\frac34,1$, and steps with
reward 0 or 1 print a loss of exactly zero because their advantages vanish. Large loss
magnitudes simply mean that the group disagreed strongly. They do not signal trouble. The
quantities to watch are the \emph{moving averages} of reward and response length and,
periodically, accuracy on held-out data. On a small base model such as Qwen3-0.6B with MATH, you should expect
a few dozen GRPO steps with 4--8 rollouts of $\le512$ tokens to raise MATH-500 accuracy
substantially, moving it towards the official reasoning variant of the same model; measure it
yourself with the Lab~27 \code{run.py} on a GPU. The model mostly learns to finish its reasoning with a boxed
answer inside the budget, and secondarily to reason better.

\Cref{fig:grpo:curves} shows the same dynamics on our toy task with the batched trainer
(\code{run.py toy-batched}): format is learned almost immediately, correctness keeps
improving, and entropy decays as the policy sharpens. The greedy accuracy \emph{without}
the box requirement also rose, from 68\% after SFT to 96\%. The SFT demonstrations were
right only 60\% of the time, so RL has turned an inconsistent arithmetic skill into a
reliable one. This ``sharpening'' of behaviour the base model can already produce is
widely regarded as the main effect of RLVR at small scale.

\begin{figure}[t]
\centering
\begin{tikzpicture}
\begin{groupplot}[group style={group size=2 by 2,horizontal sep=1.6cm,vertical sep=1.2cm},
  width=0.48\linewidth,height=4.2cm,xmin=0,xmax=200,grid=major,grid style={gray!15},
  tick label style={font=\scriptsize},label style={font=\small},ylabel shift=-3pt]
\nextgroupplot[ylabel={reward (train)},ymin=0,ymax=1.05]
\addplot[sysblue!35,thin] coordinates {(1,0.141) (2,0.188) (3,0.266) (4,0.266) (5,0.250) (6,0.234) (7,0.422) (8,0.406) (9,0.422) (10,0.625) (11,0.422) (12,0.656) (13,0.344) (14,0.453) (15,0.672) (16,0.484) (17,0.594) (18,0.531) (19,0.281) (20,0.547) (21,0.547) (22,0.500) (23,0.438) (24,0.562) (25,0.344) (26,0.391) (27,0.562) (28,0.578) (29,0.688) (30,0.500) (31,0.547) (32,0.438) (33,0.609) (34,0.531) (35,0.500) (36,0.703) (37,0.781) (38,0.688) (39,0.578) (40,0.438) (41,0.516) (42,0.484) (43,0.672) (44,0.453) (45,0.750) (46,0.516) (47,0.500) (48,0.609) (49,0.672) (50,0.750) (51,0.547) (52,0.594) (53,0.656) (54,0.516) (55,0.625) (56,0.484) (57,0.672) (58,0.641) (59,0.594) (60,0.781) (61,0.547) (62,0.703) (63,0.750) (64,0.453) (65,0.594) (66,0.672) (67,0.812) (68,0.562) (69,0.750) (70,0.688) (71,0.641) (72,0.828) (73,0.578) (74,0.703) (75,0.812) (76,0.469) (77,0.875) (78,0.609) (79,0.609) (80,0.750) (81,0.812) (82,0.734) (83,0.703) (84,0.828) (85,0.672) (86,0.750) (87,0.781) (88,0.875) (89,0.688) (90,0.641) (91,0.766) (92,0.891) (93,0.859) (94,0.781) (95,0.578) (96,0.703) (97,0.750) (98,0.656) (99,0.750) (100,0.875) (101,0.625) (102,0.797) (103,0.844) (104,0.594) (105,0.750) (106,0.828) (107,0.812) (108,0.766) (109,0.922) (110,0.750) (111,0.891) (112,0.781) (113,0.875) (114,0.875) (115,0.828) (116,0.938) (117,0.844) (118,0.906) (119,0.938) (120,0.672) (121,0.797) (122,0.891) (123,0.984) (124,0.844) (125,0.594) (126,0.797) (127,0.875) (128,0.781) (129,0.875) (130,0.766) (131,0.891) (132,0.938) (133,0.672) (134,0.906) (135,0.984) (136,0.688) (137,0.906) (138,0.938) (139,0.781) (140,0.906) (141,0.625) (142,0.812) (143,0.906) (144,0.953) (145,0.812) (146,0.812) (147,0.688) (148,0.750) (149,0.922) (150,0.875) (151,0.734) (152,0.938) (153,0.812) (154,0.875) (155,0.719) (156,0.719) (157,0.844) (158,0.859) (159,0.781) (160,0.984) (161,0.859) (162,0.750) (163,0.938) (164,0.953) (165,0.938) (166,0.953) (167,0.828) (168,0.922) (169,0.625) (170,0.688) (171,0.641) (172,0.828) (173,0.859) (174,0.781) (175,0.875) (176,0.844) (177,0.812) (178,0.953) (179,0.938) (180,0.828) (181,0.859) (182,0.812) (183,0.953) (184,0.922) (185,0.797) (186,0.969) (187,0.734) (188,0.672) (189,0.969) (190,0.953) (191,0.953) (192,0.812) (193,0.797) (194,0.906) (195,0.891) (196,0.703) (197,0.859) (198,0.938) (199,0.938) (200,0.906)};
\addplot[sysblue,very thick] coordinates {(1,0.141) (3,0.198) (5,0.222) (7,0.252) (9,0.288) (11,0.331) (13,0.357) (15,0.384) (17,0.403) (19,0.403) (21,0.430) (23,0.455) (25,0.474) (27,0.489) (29,0.511) (31,0.511) (33,0.513) (35,0.509) (37,0.529) (39,0.552) (41,0.545) (43,0.555) (45,0.570) (47,0.573) (49,0.574) (51,0.587) (53,0.597) (55,0.602) (57,0.586) (59,0.584) (61,0.603) (63,0.618) (65,0.610) (67,0.634) (69,0.635) (71,0.637) (73,0.645) (75,0.663) (77,0.673) (79,0.672) (81,0.684) (83,0.683) (85,0.705) (87,0.708) (89,0.720) (91,0.724) (93,0.741) (95,0.734) (97,0.739) (99,0.748) (101,0.745) (103,0.755) (105,0.748) (107,0.753) (109,0.759) (111,0.771) (113,0.766) (115,0.784) (117,0.800) (119,0.822) (121,0.820) (123,0.832) (125,0.837) (127,0.838) (129,0.837) (131,0.838) (133,0.835) (135,0.845) (137,0.835) (139,0.829) (141,0.832) (143,0.824) (145,0.841) (147,0.832) (149,0.833) (151,0.830) (153,0.838) (155,0.823) (157,0.821) (159,0.817) (161,0.833) (163,0.831) (165,0.838) (167,0.852) (169,0.845) (171,0.831) (173,0.828) (175,0.831) (177,0.836) (179,0.848) (181,0.841) (183,0.845) (185,0.836) (187,0.832) (189,0.837) (191,0.866) (193,0.862) (195,0.869) (197,0.864) (199,0.863)};
\nextgroupplot[ylabel={greedy accuracy},ymin=0,ymax=1.05]
\addplot[keygreen,very thick,mark=*,mark size=1.2pt] coordinates {(0,0.050) (10,0.710) (20,0.610) (30,0.610) (40,0.680) (50,0.780) (60,0.770) (70,0.800) (80,0.840) (90,0.860) (100,0.870) (110,0.880) (120,0.860) (130,0.900) (140,0.880) (150,0.890) (160,0.930) (170,0.890) (180,0.940) (190,0.900) (200,0.960)};
\nextgroupplot[ylabel={response length},ymin=3.5,ymax=7,xlabel={GRPO step}]
\addplot[sysblue!35,thin] coordinates {(1,4.797) (2,4.625) (3,5.297) (4,5.031) (5,5.844) (6,6.000) (7,6.078) (8,5.578) (9,6.391) (10,6.016) (11,5.547) (12,6.141) (13,6.094) (14,6.203) (15,5.203) (16,5.969) (17,5.797) (18,6.266) (19,5.219) (20,5.859) (21,5.594) (22,6.078) (23,6.516) (24,5.984) (25,5.500) (26,5.344) (27,5.484) (28,5.297) (29,6.938) (30,5.047) (31,5.469) (32,5.531) (33,6.031) (34,6.234) (35,5.078) (36,5.984) (37,5.297) (38,5.047) (39,5.781) (40,4.703) (41,5.312) (42,5.312) (43,5.219) (44,5.625) (45,5.734) (46,4.828) (47,4.938) (48,4.734) (49,5.094) (50,5.078) (51,5.953) (52,5.266) (53,5.203) (54,4.594) (55,4.297) (56,5.297) (57,4.766) (58,4.547) (59,4.484) (60,5.406) (61,4.859) (62,4.703) (63,4.547) (64,4.812) (65,4.922) (66,4.719) (67,4.672) (68,4.750) (69,5.234) (70,4.625) (71,4.641) (72,5.438) (73,4.781) (74,4.406) (75,4.750) (76,4.516) (77,4.641) (78,6.422) (79,5.031) (80,4.562) (81,4.531) (82,5.125) (83,4.391) (84,4.609) (85,4.641) (86,4.750) (87,4.828) (88,4.125) (89,4.219) (90,4.500) (91,4.484) (92,4.500) (93,4.484) (94,4.516) (95,4.359) (96,4.484) (97,4.719) (98,4.500) (99,4.609) (100,4.578) (101,4.891) (102,4.500) (103,4.594) (104,4.625) (105,4.641) (106,4.688) (107,4.500) (108,4.250) (109,4.484) (110,4.156) (111,4.391) (112,4.453) (113,4.500) (114,4.609) (115,4.266) (116,4.219) (117,4.469) (118,4.516) (119,4.250) (120,4.375) (121,4.750) (122,4.750) (123,4.484) (124,4.344) (125,4.891) (126,4.734) (127,4.609) (128,4.250) (129,4.500) (130,4.500) (131,4.375) (132,4.484) (133,4.625) (134,4.516) (135,4.609) (136,4.625) (137,4.266) (138,4.125) (139,4.500) (140,4.406) (141,4.766) (142,4.500) (143,4.250) (144,4.250) (145,4.250) (146,4.344) (147,4.578) (148,4.438) (149,4.250) (150,4.750) (151,4.344) (152,4.516) (153,4.641) (154,4.625) (155,4.375) (156,4.609) (157,4.391) (158,4.438) (159,4.734) (160,4.500) (161,4.375) (162,4.500) (163,4.641) (164,4.500) (165,4.500) (166,4.250) (167,4.953) (168,4.609) (169,4.625) (170,4.406) (171,4.234) (172,4.500) (173,4.609) (174,4.719) (175,4.375) (176,4.641) (177,4.250) (178,4.625) (179,4.516) (180,4.000) (181,4.672) (182,4.344) (183,4.375) (184,4.625) (185,4.484) (186,4.625) (187,4.500) (188,4.359) (189,4.375) (190,4.125) (191,4.250) (192,4.375) (193,4.266) (194,4.609) (195,4.469) (196,4.375) (197,4.625) (198,4.500) (199,4.500) (200,4.750)};
\addplot[sysblue,very thick] coordinates {(1,4.797) (3,4.906) (5,5.119) (7,5.382) (9,5.516) (11,5.564) (13,5.649) (15,5.656) (17,5.683) (19,5.689) (21,5.737) (23,5.871) (25,5.902) (27,5.839) (29,5.852) (31,5.800) (33,5.766) (35,5.762) (37,5.737) (39,5.705) (41,5.633) (43,5.530) (45,5.523) (47,5.470) (49,5.350) (51,5.376) (53,5.321) (55,5.200) (57,5.139) (59,5.049) (61,5.062) (63,4.998) (65,4.916) (67,4.898) (69,4.905) (71,4.817) (73,4.805) (75,4.818) (77,4.773) (79,4.894) (81,4.835) (83,4.848) (85,4.824) (87,4.834) (89,4.752) (91,4.737) (93,4.676) (95,4.662) (97,4.664) (99,4.547) (101,4.566) (103,4.545) (105,4.545) (107,4.526) (109,4.545) (111,4.523) (113,4.522) (115,4.522) (117,4.496) (119,4.479) (121,4.462) (123,4.469) (125,4.467) (127,4.475) (129,4.476) (131,4.492) (133,4.500) (135,4.513) (137,4.523) (139,4.516) (141,4.518) (143,4.494) (145,4.457) (147,4.436) (149,4.433) (151,4.444) (153,4.446) (155,4.440) (157,4.445) (159,4.473) (161,4.458) (163,4.477) (165,4.502) (167,4.516) (169,4.544) (171,4.521) (173,4.519) (175,4.523) (177,4.518) (179,4.516) (181,4.506) (183,4.485) (185,4.491) (187,4.487) (189,4.462) (191,4.448) (193,4.425) (195,4.424) (197,4.430) (199,4.423)};
\nextgroupplot[ylabel={entropy (nats)},ymin=0,ymax=0.7,xlabel={GRPO step}]
\addplot[sysblue!35,thin] coordinates {(1,0.527) (2,0.499) (3,0.478) (4,0.480) (5,0.436) (6,0.425) (7,0.422) (8,0.421) (9,0.353) (10,0.342) (11,0.412) (12,0.357) (13,0.439) (14,0.303) (15,0.372) (16,0.308) (17,0.347) (18,0.346) (19,0.369) (20,0.320) (21,0.312) (22,0.303) (23,0.334) (24,0.317) (25,0.286) (26,0.296) (27,0.297) (28,0.288) (29,0.275) (30,0.304) (31,0.280) (32,0.197) (33,0.224) (34,0.203) (35,0.293) (36,0.214) (37,0.194) (38,0.263) (39,0.276) (40,0.298) (41,0.224) (42,0.237) (43,0.235) (44,0.280) (45,0.224) (46,0.274) (47,0.298) (48,0.278) (49,0.184) (50,0.196) (51,0.197) (52,0.239) (53,0.172) (54,0.249) (55,0.261) (56,0.237) (57,0.163) (58,0.223) (59,0.250) (60,0.175) (61,0.169) (62,0.186) (63,0.226) (64,0.216) (65,0.229) (66,0.155) (67,0.126) (68,0.226) (69,0.157) (70,0.156) (71,0.198) (72,0.173) (73,0.119) (74,0.214) (75,0.086) (76,0.224) (77,0.146) (78,0.194) (79,0.159) (80,0.150) (81,0.171) (82,0.160) (83,0.149) (84,0.132) (85,0.206) (86,0.151) (87,0.102) (88,0.167) (89,0.161) (90,0.118) (91,0.157) (92,0.089) (93,0.080) (94,0.135) (95,0.145) (96,0.179) (97,0.128) (98,0.161) (99,0.139) (100,0.094) (101,0.144) (102,0.115) (103,0.132) (104,0.157) (105,0.125) (106,0.111) (107,0.110) (108,0.136) (109,0.104) (110,0.159) (111,0.174) (112,0.165) (113,0.129) (114,0.102) (115,0.157) (116,0.086) (117,0.127) (118,0.091) (119,0.100) (120,0.125) (121,0.080) (122,0.073) (123,0.064) (124,0.111) (125,0.169) (126,0.105) (127,0.099) (128,0.093) (129,0.116) (130,0.082) (131,0.087) (132,0.094) (133,0.144) (134,0.105) (135,0.077) (136,0.123) (137,0.079) (138,0.085) (139,0.104) (140,0.104) (141,0.132) (142,0.084) (143,0.070) (144,0.035) (145,0.064) (146,0.065) (147,0.107) (148,0.085) (149,0.078) (150,0.091) (151,0.090) (152,0.052) (153,0.082) (154,0.085) (155,0.084) (156,0.052) (157,0.071) (158,0.063) (159,0.068) (160,0.054) (161,0.078) (162,0.090) (163,0.075) (164,0.067) (165,0.078) (166,0.073) (167,0.098) (168,0.051) (169,0.109) (170,0.103) (171,0.116) (172,0.120) (173,0.087) (174,0.073) (175,0.056) (176,0.063) (177,0.092) (178,0.058) (179,0.067) (180,0.075) (181,0.101) (182,0.103) (183,0.047) (184,0.065) (185,0.076) (186,0.060) (187,0.050) (188,0.093) (189,0.036) (190,0.049) (191,0.065) (192,0.072) (193,0.080) (194,0.056) (195,0.045) (196,0.119) (197,0.054) (198,0.079) (199,0.048) (200,0.051)};
\addplot[sysblue,very thick] coordinates {(1,0.527) (3,0.501) (5,0.484) (7,0.467) (9,0.449) (11,0.436) (13,0.430) (15,0.418) (17,0.407) (19,0.402) (21,0.387) (23,0.370) (25,0.354) (27,0.342) (29,0.331) (31,0.323) (33,0.304) (35,0.295) (37,0.283) (39,0.274) (41,0.268) (43,0.260) (45,0.255) (47,0.254) (49,0.249) (51,0.239) (53,0.239) (55,0.240) (57,0.239) (59,0.236) (61,0.227) (63,0.224) (65,0.221) (67,0.206) (69,0.203) (71,0.201) (73,0.195) (75,0.184) (77,0.183) (79,0.177) (81,0.175) (83,0.170) (85,0.165) (87,0.164) (89,0.161) (91,0.157) (93,0.151) (95,0.150) (97,0.147) (99,0.144) (101,0.140) (103,0.137) (105,0.134) (107,0.132) (109,0.128) (111,0.131) (113,0.137) (115,0.136) (117,0.131) (119,0.126) (121,0.124) (123,0.119) (125,0.119) (127,0.118) (129,0.116) (131,0.108) (133,0.105) (135,0.101) (137,0.101) (139,0.101) (141,0.102) (143,0.103) (145,0.094) (147,0.093) (149,0.090) (151,0.091) (153,0.086) (155,0.085) (157,0.081) (159,0.078) (161,0.073) (163,0.074) (165,0.076) (167,0.076) (169,0.076) (171,0.078) (173,0.081) (175,0.079) (177,0.081) (179,0.080) (181,0.083) (183,0.082) (185,0.082) (187,0.079) (189,0.077) (191,0.072) (193,0.069) (195,0.068) (197,0.069) (199,0.069)};
\end{groupplot}
\end{tikzpicture}
\caption{A GRPO run on the toy addition task (8 prompts $\times$ 8 rollouts per step,
learning rate $3\times10^{-4}$, no KL, no clipping), measured on a CPU in this book's
build environment. Light lines are per-step values, dark lines a 20-step moving average.
Greedy accuracy (a correct \emph{and} boxed answer, over all 100 prompts) rises from 5\% to
96\%. The run learns the format within about ten steps and then keeps sharpening
correctness. Entropy falls steadily as the policy becomes more deterministic.}
\label{fig:grpo:curves}
\end{figure}

\begin{systemslens}[where a GRPO step spends its time]
Over the 200 steps of \cref{fig:grpo:curves} (64 rollouts per step, measured on one CPU core)
the trainer logged 26.8\,s sampling, 9.0\,s computing old log-probs and entropies under
\code{no\_grad}, and 28.4\,s in differentiable forward, backward and optimiser updates.
Sampling was \emph{only} about 40\% of the step here because toy completions are about five tokens
long. On a GPU with realistic reasoning lengths the balance flips. Using the roofline numbers
from \cref{sec:grpo:rollouts}, for Qwen3-0.6B with 64 rollouts of 1024 tokens:

\smallskip
\centering
\begin{tikzpicture}
\begin{axis}[xbar stacked,width=0.9\linewidth,bar width=9pt,
  xlabel={seconds per GRPO step},symbolic y coords={CPU toy measured,H100 lower bound},ytick=data,
  y tick label style={font=\scriptsize},x tick label style={font=\scriptsize},
  legend style={font=\scriptsize,at={(0.5,1.05)},anchor=south,legend columns=3,draw=none},xmin=0,height=3.6cm,
  enlarge y limits=0.5]
\addplot[fill=fcompute!60,draw=none] coordinates {(0.134,CPU toy measured) (2.09,H100 lower bound)};
\addplot[fill=fmemory!60,draw=none] coordinates {(0.045,CPU toy measured) (0.25,H100 lower bound)};
\addplot[fill=fcomm!60,draw=none] coordinates {(0.142,CPU toy measured) (0.74,H100 lower bound)};
\legend{sampling,old log-probs,fwd+bwd+step}
\end{axis}
\end{tikzpicture}
\raggedright

\smallskip
The H100 bars are lower bounds: the bandwidth roofline for decoding, and 40\% MFU
for the training passes. Real samplers sit well above their bound, and reasoning
rollouts are far longer than 1024 tokens, so in practice sampling usually accounts for most of
the wall-clock time. The toy run is the exception that proves the rule: shrink the
completions and the training side dominates again. That is why the balance has to be
\emph{measured} per workload, which Lab~27 asks you to do with \code{time\_breakdown}.
\end{systemslens}

\paragraph{Batching and multiple GPUs.} The loop above handles one prompt per step and
generates its $G$ rollouts sequentially, the clearest form and the slowest. The library trainer
batches all $P\times G$ rollouts of a step into one decode batch, computes log-probs for all
of them in one padded forward pass, and is the version used for every experiment in
\cref{ch:grpo2}. Batched RL rollouts are revisited in \cref{ch:serving} (padding and masks),
and \cref{ch:rlsys} distributes the sampler and trainer across GPUs.

% ======================================================================
\section{Loading and evaluating checkpoints}\label{sec:grpo:ckpt}

The loop saves \code{model.state\_dict()} every few steps and on keyboard interrupt, so a
long run can be stopped at any point without losing work. Restoring is the usual one-liner,
\code{model.load\_state\_dict(torch.load(path))}, and the restored model plugs straight into
the evaluation harness of \cref{ch:eval}:

\begin{lstlisting}[style=shell,caption={Evaluating an RL checkpoint on MATH-500 with the Chapter 21 harness.},label={lst:grpo:eval}]
python labs/lab21_eval/run.py --dataset math500 --model-dir Qwen3-0.6B-Base \
    --checkpoint out17/ckpt/grpo-step00050.pt --max-new-tokens 2048
\end{lstlisting}

Two evaluation details matter. First, evaluate with a \emph{larger} token budget than the
training budget. A model trained with 512-token rollouts has learned to box its answer
within 512 tokens, and that should not cap its test-time reasoning. Second, track
average response length alongside accuracy. In DeepSeek-R1's training, response length
grew steadily as the model learned to reason longer and check its work. A tight training
budget suppresses that growth, which is a direct compute--quality trade-off: longer
rollouts cost linearly more decode steps and quadratically more attention, and the KV cache
grows with them.

On the toy task, \code{evaluate\_accuracy} decodes greedily on all 100 prompts. The SFT
start scores 5\% (correct \emph{and} boxed) and 68\% (correct, any format). After 200 GRPO
steps both are 96\%. More steps are not automatically better. Plain GRPO without clipping
or KL can become unstable in long runs, as the reward collapses and gibberish appears in
\cref{ch:grpo2}. Keep several checkpoints and select by validation accuracy, not by the
last step.

\begin{napkinmath}[checkpoint budget]
A Qwen3-0.6B \code{state\_dict} in bf16 is $0.6\times10^9\times2\,\text{B}\approx1.2$\,GB
(twice that if it was trained in fp32). Saving every 5 steps for
500 steps produces 100 checkpoints, or 120\,GB. Saving every 50 steps produces 12\,GB. An 8B
policy with optimiser state (needed to \emph{resume}, not just evaluate) is
$8\times10^9\times16\,\text{B}=128$\,GB per checkpoint, and writing it to network storage at
2\,GB/s stalls a synchronous trainer for over a minute. Hence the asynchronous, sharded
checkpointing of \cref{ch:pretrain}.
\end{napkinmath}

% ======================================================================
\section{Where the memory goes: a GRPO step at scale}\label{sec:grpo:systems}

We close the chapter by assembling the systems costs scattered through it into one budget,
which is the starting point for \cref{ch:rlsys}. Consider one GRPO step on an 8B policy with
$P=32$ prompts, $G=8$ rollouts each and up to 4096 new tokens. \Cref{tab:grpo:budget} lists
what must exist in GPU memory and when.

\begin{table}[ht]
\centering\small
\begin{tabularx}{\linewidth}{@{}lXr@{}}
\toprule
Phase & What lives in memory & Size (8B, bf16)\\
\midrule
sampling & policy weights + KV cache for 256 sequences & 16\,GB + $\approx$\,150\,GB$^\dagger$\\
old log-probs & weights + one micro-batch of activations (no grad) & 16\,GB + small\\
reference log-probs (if KL) & a second, frozen copy of the weights & +16\,GB\\
training & weights, grads, AdamW states, fp32 master & 128\,GB\\
 & activations of one micro-batch (checkpointed) & 5--20\,GB\\
 & $(T\times V)$ logits for the log-softmax & $\approx$\,5\,GB per 8k tokens\\
\bottomrule
\end{tabularx}
\caption{Memory of one GRPO step for an 8B dense policy. $^\dagger$KV cache:
$2\cdot L\cdot H_{kv}\cdot d_h\cdot\text{bytes}=2\cdot36\cdot8\cdot128\cdot2=147$\,KB per
token, times $256\times4096$ tokens $\approx154$\,GB. No single GPU holds this, so sampling
must be sharded or run in waves.}
\label{tab:grpo:budget}
\end{table}

Three observations drive the design of every RL framework:
\begin{enumerate}
\item \textbf{The two phases want different parallelism.} Sampling wants many replicas
with lots of free memory for KV cache (tensor parallel at most within a node). Training wants the
optimiser state sharded (FSDP/ZeRO-3, \cref{ch:distributed}). Running both on the same GPUs
(colocation) means resharding the weights between layouts every step. Running them on
separate GPUs (disaggregation) means shipping 16\,GB of new weights to the samplers every
step.
\item \textbf{The sampler and trainer should not wait for each other.} Synchronous GRPO
leaves trainer GPUs idle during sampling and sampler GPUs idle during training. Overlapping
them makes the data slightly off-policy (generated by an older $\theta$), which is
exactly what the importance ratios and clipping of \cref{ch:grpo2} correct for.
\item \textbf{Variance in response length is a scheduling problem.} The step ends when the
\emph{longest} rollout ends. With reasoning lengths spanning 100 to \num{30000} tokens, most
of the batch sits finished while a few stragglers decode at batch size 1, the least
efficient point on the roofline.
\end{enumerate}

% ======================================================================
\section{Summary}\label{sec:grpo:summary}
\begin{itemize}
\item RL post-training optimises a property of whole outputs. RLHF learns a reward model
from human rankings and optimises it with PPO. RLVR replaces the reward model by a
deterministic verifier (answer checking, unit tests).
\item The policy-gradient theorem turns a non-differentiable reward into a weighting of
$\nabla\log\pitheta(y\mid x)$. Subtracting a baseline keeps the estimator unbiased and cuts its
variance.
\item GRPO's baseline is the mean reward of $G$ rollouts of the same prompt, and its advantages
are those rewards standardised within the group. There is no value model, so memory is lower,
at the price of $G\times$ more sampling.
\item Groups whose rollouts all receive the same reward contribute zero gradient. With binary
rewards that happens with probability $p^G+(1-p)^G$.
\item The loss implemented here is $-\frac1G\sum_iA_i\log\pitheta(y_i\mid x)$ with
\emph{summed} token log-probs. Its value is not a quality measure; its gradient is what
matters.
\item Rollout code must use \code{no\_grad} (not \code{inference\_mode}), keep the EOS token,
and return ids rather than text. Log-probs are recomputed in a separate differentiable
pass, and that $(T\times V)$ log-softmax is a major memory consumer.
\item The training loop is otherwise standard. Track moving-average reward, response
length and held-out accuracy, and checkpoint often.
\item At scale a GRPO step is an inference job glued to a training job with different memory
and parallelism needs, the central problem of \cref{ch:rlsys}.
\end{itemize}

% ======================================================================
\section{Exercises}\label{sec:grpo:exercises}

\begin{exercise}[format-aware reward shaping]
Implement \code{reward\_partial}, which returns 1.0 for a boxed correct answer, 0.5 for a
correct but unboxed final number, and 0 otherwise. For the four rollouts of
\cref{fig:grpo:step}, recompute the advantages and the loss. Which rollout's push changes
sign? Argue whether partial credit speeds up or slows down learning of the format.
\end{exercise}

\begin{exercise}[zero-advantage groups]
Force all rewards in a group to 0 (then to 1). Verify that the loss \emph{and every
parameter gradient} are exactly zero. Explain why this is the desired behaviour, and
compute the expected fraction of wasted rollouts per step for a dataset whose per-prompt
success rates are uniform on $[0,1]$, for $G=4$ and $G=16$.
\end{exercise}

\begin{exercise}[binary-reward advantages]
Derive $A^+=(1-p)/s$ and $A^-=-p/s$ with $s=\sqrt{\frac{G}{G-1}p(1-p)}$, and show that
$\sum_iA_i=0$ and $\frac1{G-1}\sum_iA_i^2\approx1$ for every $0<p<1$. What happens to
$\max|A_i|$ as $p\to 1/G$ for large $G$?
\end{exercise}

\begin{exercise}[the baseline's bias]
GRPO's baseline includes $r_i$ itself. Show that the leave-one-out baseline
$b_i=\frac1{G-1}\sum_{j\ne i}r_j$ (RLOO) yields an unbiased gradient, and that
$r_i-b_i=\frac{G}{G-1}(r_i-\bar r)$. Conclude that, before std scaling, GRPO's mean-centred
advantage is RLOO up to a constant factor.
\end{exercise}

\begin{exercise}[sum versus mean]
Write the per-token gradient weight of a rollout of length $|y|$ under (a) the summed loss
of this chapter and (b) a per-sequence mean. For two correct answers of lengths 20 and 2000,
compute the ratio of per-token pushes under each scheme. Which scheme pushes the policy
towards shorter correct answers, and why?
\end{exercise}

\begin{exercise}[the \code{inference\_mode} trap]
Decorate \code{sample\_response} with \code{@torch.}\allowbreak\code{inference\_mode()} instead of
\code{@torch.no\_grad()} and run \code{compute\_grpo\_loss}. Record the exact error,
explain which tensor triggers it, and find the smallest change that fixes it without
giving up inference mode for sampling.
\end{exercise}

\begin{exercise}[logits memory]
Compute the size of the float32 $(N,T,|\mathcal V|)$ log-softmax for $N=16$,
$T=4096$ and Qwen3's vocabulary, with and without its gradient. Implement a chunked
variant of \code{sequence\_logprob} (process $C$ positions at a time). Measure peak memory
with \code{torch.cuda.max\_memory\_allocated} (or \code{tracemalloc} on CPU) for several $C$.
\end{exercise}

\begin{exercise}[measure the step]
Use \code{time\_breakdown} on the toy task for $G\in\{2,4,8,16\}$ and
\code{max\_new\_tokens}$\in\{4,16,64\}$ (force long outputs by masking EOS). Fit
$t_{\text{sample}}\approx a\cdot G\cdot T$ and $t_{\text{train}}\approx b\cdot G\cdot T$. Then
replace the sequential sampler by \code{sample\_rollouts} (one batch) and explain the change
in $a$ using the roofline from \cref{ch:inference}.
\end{exercise}

\begin{exercise}[napkin: a 4B run]
For Qwen3-4B with $P=16$, $G=16$, 256-token prompts and 4096 new tokens, compute (i) the KV
cache at the end of sampling, (ii) the bandwidth-bound sampling time on one H100, and (iii) the
training FLOPs and their time at 40\% MFU. Does the KV cache fit next to the weights on one
80\,GB GPU? If not, how many sampling waves are needed?
\end{exercise}

\begin{exercise}[break the verifier]
Write five completions that a naive verifier (first \verb|\boxed{}|, substring
match, no normalisation) scores incorrectly, in both directions. Harden
\code{reward\_rlvr} against them and add them as unit tests. Which loopholes would an RL
policy be most likely to discover, and why?
\end{exercise}

\begin{exercise}[held-out generalisation]
Rebuild the toy with \code{make\_examples(holdout=20)}: SFT and GRPO on 80 facts and
evaluate on the 20 held-out ones. Does GRPO improve held-out \emph{format}? Held-out
\emph{correctness}? Relate your findings to the separation between MATH and MATH-500.
\end{exercise}

% ======================================================================
\section{Lab 27: GRPO from scratch}\label{sec:grpo:lab}

\begin{labbox}{27}{GRPO from scratch}
\labdir{lab27_grpo}
\labsection{Goals}
Implement every stage of RLVR with GRPO, see a policy learn a verifiable task on a CPU, and
measure where a GRPO step spends its time.

\labsection{Setup}
\begin{lstlisting}[style=shell]
cd code
LAB_IMPL=starter  python -m pytest -q labs/lab27_grpo    # your implementation
LAB_IMPL=solution python -m pytest -q labs/lab27_grpo    # reference
\end{lstlisting}
The toy task, tokenizer and SFT warm start are provided in \code{llmsys.rl.toy}; the tiny
Qwen3-shaped model comes from \code{llmsys.qwen3}.

\labsection{Tasks} (each is a function in \code{starter.py})
\begin{enumerate}
\item \code{sample\_response}: KV-cached sampler with temperature and top-$p$ that keeps EOS and returns
ids, prompt length and text (\cref{lst:grpo:sample}).
\item \code{reward\_rlvr} and \code{reward\_partial}: strict and partial-credit rewards.
\item \code{compute\_advantages}: group standardisation with unbiased std and $\epsilon$.
\item \code{sequence\_logprob}: differentiable summed log-probability with the correct
off-by-one (\cref{fig:grpo:shift}).
\item \code{policy\_gradient\_loss}: $-\operatorname{mean}(A\cdot\log\pi)$ with detached advantages.
\item \code{compute\_grpo\_loss}: stages 1--5 for one example.
\item \code{train\_rlvr\_grpo}: the loop, with clipping, CSV logging and checkpoints.
\item \code{time\_breakdown}: seconds spent sampling, in the log-prob pass, and in backward.
\end{enumerate}

\labsection{Measurements}
\begin{enumerate}
\item \code{python labs/lab17\_grpo/run.py toy --impl starter --steps 150}: your
sequential implementation should lift the moving-average reward from about 0.1--0.2 to above 0.7.
\item \code{run.py toy-batched --steps 200}: reproduce \cref{fig:grpo:curves}. Plot
reward, greedy accuracy, length and entropy with \code{llmsys.rl.metrics.plot\_metrics}.
\item Compare wall-clock per step of the sequential and batched versions and report the
speed-up of sampling alone. Expect it to grow with $G$.
\item Fill in a table of \code{time\_breakdown} for $G\in\{2,4,8\}$ and explain the
trends (Exercise~27.8).
\end{enumerate}

\labsection{Deliverables}
Passing tests; a plot of your training curves; a one-paragraph explanation of the
time breakdown; the greedy accuracy before and after GRPO.

\labsection{Stretch goals}
\begin{itemize}
\item On a GPU, run \code{run.py qwen3} on Qwen3-0.6B-Base with MATH for 50 steps, then
evaluate the checkpoint on MATH-500 with Lab~21's harness.
\item Replace the group-mean baseline by the leave-one-out baseline (Exercise~27.4) and
compare learning curves over three seeds.
\item Make sampling and training overlap: generate step $k{+}1$'s rollouts in a background
thread while step $k$ trains, and measure the speed-up and the (small) policy lag. This
previews \cref{ch:rlsys}.
\end{itemize}
\end{labbox}
